\section{Chapter Preface}
A robust abstraction should reveal sameness across domains that appear unrelated. Applications in this chapter are therefore not demonstrations of breadth alone; they are tests of structural invariance.

Computer vision, robotics, physiology, program synthesis, and distributed systems each become instances of a single reconstruction architecture once their latent states, observations, and constraints are made explicit.

The chapter emphasizes representation theorems over domain anecdotes, showing how one formal tuple can unify heterogeneous practice.

\section{Derivations and Formal Results}
Applications are presented as instances of a common CLIO tuple \((\stateSpace,\observationSpace,\projection,\constraintSet,F,J)\).

\begin{description}
  \item[Computer vision] Inverse rendering model: \(y=\mathcal R(x)+\eta\).
  \item[Robotics] State estimation as admissible latent reconstruction under dynamics and sensor constraints.
  \item[Biology] Physiological admissibility regions plus repair operators for homeostatic recovery.
  \item[Program synthesis] Programs as latent structures constrained by behavioral specifications.
  \item[Distributed systems] Global consistency from local propagation over shared constraints.
\end{description}

\begin{theorem}[Representation principle]
Each domain above can be encoded as a CLIO system by choosing domain-specific state variables, observation map, and admissibility constraints while preserving common reconstruction dynamics.
\end{theorem}
